% test-kksref.tex
% \kksref（基準ボックス）機構と中央寄せの検証用。
% コンパイル: latexmk -r .latexmkrc test-kksref.tex   (LuaLaTeX 専用)
%
% KKran があれば \Rrnum を使い、無ければ \textit{...} でローマ数字小文字を代用する。
\documentclass{jlreq}
\usepackage[tsumesuji=1]{KKsymbols}
\usepackage{luatexja-otf}

% --- KKran(\Rrnum) フォールバック ----------------------------------------
\IfFileExists{KKran.sty}{%
  \usepackage{KKran}%
}{%
  % KKran が無い環境用の簡易代用。lowercase roman numerals。
  \makeatletter
  \newcommand{\Rrnum}[1]{\textit{\romannumeral #1}}%
  \makeatother
}

% --- Hiragino（環境に無ければ適宜差し替え） --------------------------------
\makeatletter
\RequirePackage[no-math]{fontspec}
\RequirePackage[no-math,match,scale=1]{luatexja-fontspec}
\RequirePackage[hiragino-pro,deluxe,expert]{luatexja-preset}
\makeatother

\fboxsep=0pt\fboxrule=.1pt
\begin{document}

\section*{1. 問題1: ローマ数字小文字の見た目揃え}

\noindent\textbf{(a) 何もしない（基準ボックス無し）— 字ごとにバラつくはず:}\par
\fbox{\maru{\Rrnum{1}}}\fbox{\maru{\Rrnum{2}}}\fbox{\maru{\Rrnum{3}}}%
\fbox{\maru{\Rrnum{6}}}\fbox{\maru{\Rrnum{7}}}\fbox{\maru{\Rrnum{8}}}

\bigskip
\noindent\textbf{(b) 従来の手書き回避策 \texttt{\textbackslash vphantom\{\textbackslash Rrnum\{6\}\}}:}\par
\fbox{\maru{\vphantom{\Rrnum{6}}\Rrnum{1}}}\fbox{\maru{\vphantom{\Rrnum{6}}\Rrnum{2}}}%
\fbox{\maru{\vphantom{\Rrnum{6}}\Rrnum{3}}}\fbox{\maru{\vphantom{\Rrnum{6}}\Rrnum{6}}}%
\fbox{\maru{\vphantom{\Rrnum{6}}\Rrnum{7}}}\fbox{\maru{\vphantom{\Rrnum{6}}\Rrnum{8}}}

\bigskip
\noindent\textbf{(c) 新機構 \texttt{\textbackslash kksref\{\textbackslash Rrnum\{6\}\}} — (b) と一致すれば成功:}\par
{\kksref{\Rrnum{6}}%
 \fbox{\maru{\Rrnum{1}}}\fbox{\maru{\Rrnum{2}}}\fbox{\maru{\Rrnum{3}}}%
 \fbox{\maru{\Rrnum{6}}}\fbox{\maru{\Rrnum{7}}}\fbox{\maru{\Rrnum{8}}}}

\bigskip
\noindent\textbf{(d) 他記号でも:}\par
{\kksref{\Rrnum{6}}%
 \fbox{\seihou{\Rrnum{3}}}\fbox{\kuroseihou{\Rrnum{8}}}\fbox{\seimaru{\Rrnum{7}}}%
 \fbox{\kakko{\Rrnum{4}}}}

\section*{2. 後方互換チェック（基準ボックス未設定 = 従来出力のまま）}

\noindent\maru{1}\maru{97}\maru{だ}\maru{ばばば}\maru*{m}\maru{Qjg}\par
\kuromaru{1}\nmaru{1}\jegg{1}\seihou{5}\seimaru{5}\hishi{3}\kakko{12}

\section*{3. 問題2: フォント差での中央寄せ}

\noindent
{\mcfamily 明朝: \maru{8}\seihou{8}\kuromaru{3}}\quad
{\gtfamily ゴシック: \maru{8}\seihou{8}\kuromaru{3}}\quad
{\mgfamily 丸ゴ: \maru{8}\seihou{8}\kuromaru{3}}

\end{document}
